\documentclass{article}

\usepackage{amsmath,amssymb}
\usepackage{xurl}
\usepackage[hidelinks]{hyperref}
\setlength{\emergencystretch}{2em}

\input{command}

\title{SIC--POVM Equality Families: the One-Dimensional Degeneracy}
\author{}
\date{2026-08-21}

\begin{document}
\maketitle
\gapnote{local-correction}{resolved}

\section{Source passage}

Let $P_1,\ldots,P_n\in\MN{d}$ be positive semidefinite matrices with $d\leq n$
and
\begin{align}
    \tr(P_i^2)&=1
    \qquad(1\leq i\leq n).
    \label{eq:wolf_sic_povm_linear_independence_unit_purity}
\end{align}
The proposition entitled ``SIC POVMs'' in Chapter~2 of
Ref.~\cite{Wolf2012Quantum} proves
\begin{align}
    \sum_{i\neq j}
        \left|\tr(P_iP_j)\right|^2
        &\geq
        \frac{n(n-d)^2}{(n-1)d^2}.
    \label{eq:wolf_sic_povm_linear_independence_overlap_bound}
\end{align}
At equality, the proof gives
\begin{align}
    P_i^2&=P_i,
    \label{eq:wolf_sic_povm_linear_independence_projector}\\
    \operatorname{rank}(P_i)&=1,
    \label{eq:wolf_sic_povm_linear_independence_rank_one}\\
    \sum_iP_i&=\frac{n}{d}\Id,
    \label{eq:wolf_sic_povm_linear_independence_tight_sum}\\
    \tr(P_iP_j)
        &=\frac{n-d}{(n-1)d}
        \qquad(i\neq j).
    \label{eq:wolf_sic_povm_linear_independence_constant_overlap}
\end{align}
The source then asserts that every equality family is linearly independent
\cite{Wolf2012Quantum}. This assertion appears at lines 816--823 of the local
Chapter~2 source file.

\section{The one-dimensional counterexample}

The linear-independence conclusion fails when $d=1$ and $n>1$. In one
dimension, positivity and
(\ref{eq:wolf_sic_povm_linear_independence_unit_purity}) force
\begin{align}
    P_i&=[1]
    \qquad(1\leq i\leq n).
    \label{eq:wolf_sic_povm_linear_independence_dimension_one_family}
\end{align}
For this family, each off-diagonal overlap equals one, so
\begin{align}
    \sum_{i\neq j}
        \left|\tr(P_iP_j)\right|^2
        &=n(n-1),
    \label{eq:wolf_sic_povm_linear_independence_dimension_one_left}\\
    \frac{n(n-d)^2}{(n-1)d^2}
        &=\frac{n(n-1)^2}{n-1}
    \label{eq:wolf_sic_povm_linear_independence_dimension_one_right_first}\\
        &=n(n-1).
    \label{eq:wolf_sic_povm_linear_independence_dimension_one_right_value}
\end{align}
Thus the family attains equality in
(\ref{eq:wolf_sic_povm_linear_independence_overlap_bound}). It is linearly
dependent whenever $n>1$ because all its members are equal.

\section{The corrected coefficient argument}

Assume that an equality family satisfies a linear relation
\begin{align}
    \sum_i c_iP_i&=0.
    \label{eq:wolf_sic_povm_linear_independence_linear_relation}
\end{align}
Taking the trace and using
(\ref{eq:wolf_sic_povm_linear_independence_projector}) and
(\ref{eq:wolf_sic_povm_linear_independence_rank_one}) gives
\begin{align}
    \sum_i c_i&=0.
    \label{eq:wolf_sic_povm_linear_independence_coefficient_sum}
\end{align}
Pairing
(\ref{eq:wolf_sic_povm_linear_independence_linear_relation}) with $P_j$ in
the trace inner product and using
(\ref{eq:wolf_sic_povm_linear_independence_constant_overlap}) yields
\begin{align}
    0
        &=c_j+
        \frac{n-d}{(n-1)d}\sum_{i\neq j}c_i
    \label{eq:wolf_sic_povm_linear_independence_pairing_first}\\
        &=
        \left(1-\frac{n-d}{(n-1)d}\right)c_j.
    \label{eq:wolf_sic_povm_linear_independence_pairing_prefactor}
\end{align}
For $n>1$, the prefactor satisfies
\begin{align}
    1-\frac{n-d}{(n-1)d}
        &=\frac{n(d-1)}{(n-1)d}.
    \label{eq:wolf_sic_povm_linear_independence_prefactor_value}
\end{align}
It vanishes exactly when $d=1$. Therefore every equality family is linearly
independent under the corrected condition $d>1$.

When $n=1$, the bound
(\ref{eq:wolf_sic_povm_linear_independence_overlap_bound}) is not stated
because its right-hand side has denominator $n-1$. The singleton family is
nevertheless linearly independent whenever its sole matrix has unit purity,
because that matrix is nonzero.

\section{Formalization status}

The overlap bound and its equality characterization are formalized in
\doclink{QICLean/Algebra/SICPOVMBound}. The declaration
\leanid{Matrix.sicPOVM_offDiag_overlap_sq_bound} proves
(\ref{eq:wolf_sic_povm_linear_independence_overlap_bound}), and
\leanid{Matrix.sicPOVM_offDiag_overlap_sq_eq_iff} proves the full equality
characterization
(\ref{eq:wolf_sic_povm_linear_independence_projector})--%
(\ref{eq:wolf_sic_povm_linear_independence_constant_overlap}).

The corrected linear-independence result is divided into two declarations.
The theorem
\leanid{Matrix.sicPOVM_linearIndependent_of_overlap_bound_eq} covers $d\geq2$,
and
\leanid{Matrix.singleton_linearIndependent_of_trace_sq_eq_one} covers the
singleton case. The specialized \leanid{SICPOVM} API then provides the
$n=d^2$ operator-basis and channel consequences.

\section{Verdict}

The overlap bound and equality characterization from
Ref.~\cite{Wolf2012Quantum} are formalized without an artificial restriction
on $d$. The final linear-independence claim is false for $d=1$ and $n>1$.
The corrected conclusion assumes $d>1$, with the singleton case handled
separately. This is a local source correction, and it is fully formalized.

\bibliographystyle{alpha}
\bibliography{references}

\end{document}
